% TOP_PROBLEM: {"problemId": "problem.the-cherlin-zilber-algebraicity-conjecture-for-simple-groups-of-finite-morley-rank", "canonicalTitle": "The Cherlin-Zilber Algebraicity Conjecture for Simple Groups of Finite Morley Rank", "releaseRank": 382, "primaryDomain": "logic_foundations_set_theory", "primaryDomainLabel": "Logic, foundations & set theory", "displayStatus": "open", "releaseStatus": "open"}
% !TeX program = lualatex
% ATTEMPT_STATUS: unresolved
\documentclass[11pt]{article}
\usepackage[margin=25mm]{geometry}
\usepackage{fontspec}
\setmainfont{Latin Modern Roman}
\usepackage{amsmath,amssymb,mathtools}
\usepackage{unicode-math}
\setmathfont{Latin Modern Math}
\usepackage{xurl,hyperref}
\hypersetup{colorlinks=true,urlcolor=blue}
\setcounter{secnumdepth}{0}
\setlength{\parindent}{0pt}
\setlength{\parskip}{5pt}
\setlength{\emergencystretch}{3em}
\begin{document}
\section*{\texorpdfstring{The Cherlin-Zilber Algebraicity Conjecture for Simple Groups of Finite Morley Rank}{The Cherlin-Zilber Algebraicity Conjecture for Simple Groups of Finite Morley Rank}}\label{382}
\subsection{Definitions and mathematical statement}
A definable set has Morley rank at least $\alpha+1$ if it contains infinitely many pairwise disjoint definable subsets of rank at least $\alpha$, with the usual transfinite limit rule; parameters are allowed in a sufficiently saturated elementary extension. A group has finite Morley rank when its underlying definable set has finite rank. The Cherlin--Zilber assertion says that every infinite simple group $G$ of finite Morley rank is abstractly isomorphic to
\[
 \mathbf G(K)
\]
for a linear algebraic group $\mathbf G$ over an algebraically closed field $K$. Here simple means no proper nontrivial normal subgroup, not merely no such definable subgroup. The conclusion is an algebraic realization of the group, rather than just similarities between its rank and algebraic dimension; no characteristic of $K$ is prescribed.
\par\smallskip\textit{Formulation and concept references:} \href{https://arxiv.org/abs/2606.18207}{[S1]}.

\subsection{Short English statement}
Must every infinite simple group of finite Morley rank come from a linear algebraic group over an algebraically closed field?
\subsection{Research attempt}
\paragraph{Scope and primary-source audit (27 September 2026).}
Simplicity here is abstract simplicity, and the group is infinite. The
ambient definability may have parameters; rank is computed in the
specified sufficiently saturated setting. We do not replace Morley
rank by the number of generators, ordinary cardinality, or the dimension
of a chosen representation that has not yet been constructed.

Tent's 2026 ICM account [S1] was inspected, including its final outlook.
Its sharply two-transitive constructions are not announced as examples
of finite Morley rank: the paper explicitly distinguishes the constructed
groups from the finite-rank candidates that remain to be found.
Thus the existence of exotic simple or sharply transitive groups alone
is not a counterexample to the selected conjecture.

The basic rank tools used below are standard facts of finite-Morley-rank
group theory, recorded in the primary lecture notes [E1,E2]: ranks are
preserved by definable bijections; finite unions take maximum rank;
rank adds across definable maps with constant-rank fibers; finite Morley
degree is additive on disjoint sets of the same rank; and a definably
connected group has Morley degree one. We explicitly use these facts,
not a claim that they already provide an algebraic-group classification.

The attempt extracts restrictions on any potential simple counterexample,
proves a centralizer obstruction, and carries out a field-reconstruction
calculation under a concrete additional permutation-group hypothesis.
The place where that hypothesis has not been obtained is stated at the
end.

\paragraph{Rank controls index but does not count group elements.}
If $H\le G$ is definable and has infinite index, its distinct cosets
are infinitely many pairwise disjoint definable sets of rank
$\operatorname{rk}(H)$. The definition of rank therefore gives
\[
 [G:H]=\infty\quad\Longrightarrow\quad
 \operatorname{rk}(H)<\operatorname{rk}(G).
 \tag{382.1}
\]
If the index is finite, $G$ is a finite union of definably bijective
copies of $H$, so the ranks agree. At the level of degree,
\[
 \deg(G)=[G:H]\deg(H).
 \tag{382.2}
\]
Rank zero definable sets are finite: an infinite set contains infinitely
many parameter-definable singletons, each of rank zero, and so has
rank at least one. Finiteness of degree is one of the standard finite-rank
facts listed above.

These observations imply the descending chain condition on definable
subgroups. In a strictly decreasing chain ranks can decrease only
finitely many times. Once rank is constant, each proper inclusion has
finite index at least two, and (382.2) makes the positive integer degree
strictly decrease. An infinite such chain is impossible.

Consequently an intersection of any family of definable subgroups is
already the intersection of finitely many of them. Choose a minimal
finite intersection using the descending chain condition; intersecting
with any further member cannot decrease it. This makes the abstract
intersection definable. It does not imply that every abstract subgroup
is definable or that every ascending chain stabilizes.

\paragraph{Connectedness is automatic for an infinite abstractly simple group.}
In fact an infinite simple group has no proper subgroup of finite index,
even without definability. If $H$ had index $m<\infty$, the coset
action would give a homomorphism $G\to\operatorname{Sym}(m)$.
Its kernel is normal. Simplicity makes that kernel either $1$ or $G$.
The former embeds the infinite group in a finite group; the latter
makes every coset fixed and forces $H=G$. Both contradict a proper
finite-index subgroup.

In particular the finite-rank group in the question is definably
connected and hence has Morley degree one by the standard connectedness
theorem. Its center is trivial: a nontrivial center would, by simplicity,
be the whole group, making it abelian. An abstractly simple abelian
group is cyclic of prime order. To check the latter assertion, any
nonidentity element generates a nontrivial normal subgroup, hence the
whole group; an infinite cyclic group or a finite cyclic group of
composite order has a proper nontrivial subgroup.

Similarly the group is perfect. Its derived subgroup is a normal
subgroup, and if it were trivial the group would be abelian; hence
$[G,G]=G$. This uses abstract simplicity and requires no assertion that
the derived subgroup is definable. In particular no nontrivial solvable
infinite group is among the candidates.

\paragraph{A conjugacy-class rank formula and its consequence.}
For $g\in G$, the conjugation map $x\mapsto xgx^{-1}$ has fibers
cosets of the definable centralizer $C_G(g)$. The constant-fiber rank
formula gives
\[
 \operatorname{rk}(g^G)=\operatorname{rk}(G)
                          -\operatorname{rk}(C_G(g)).
 \tag{382.3}
\]
If $C_G(g)$ is finite, its conjugacy class has full rank and is called
generic. In a degree-one group two disjoint definable generic sets
cannot exist, since their union would have degree at least two.
Therefore any two elements with finite centralizers must be conjugate.
This last statement concerns finite centralizers; it is not an assertion
that every two noncentral elements are conjugate.

\textbf{Proposition 382.1.} An infinite abstractly simple group of
finite Morley rank has a nonidentity element with an infinite proper
centralizer. In particular its rank cannot be one.

\emph{Proof.} Suppose every nonidentity element has finite centralizer.
By (382.3) and degree one they all lie in one conjugacy class. Each
such element has finite order, since its cyclic subgroup is contained
in its finite centralizer. Their orders are consequently all the same
integer $m>1$. For any prime $p$ dividing $m$, an element of order
$m$ has a power of order $p$; that power is nonidentity and must have
order $m$ as well. Hence $m=p$ is prime.

If $p=2$, every element is its own inverse, and
$ab=(ab)^{-1}=b^{-1}a^{-1}=ba$. The group is abelian, impossible.
If $p$ is odd, choose $g\ne1$. Since $g^{-1}$ is conjugate to $g$,
there is $h$ with $hgh^{-1}=g^{-1}$. Conjugation by $h$ therefore
acts by inversion on the cyclic subgroup generated by $g$.
But $h^p=1$, while the $p$th iterate of inversion, for odd $p$, is
still inversion and is not the identity on a group of order $p>2$.
This is a contradiction. Thus an infinite centralizer exists.
It is proper because the center of $G$ is trivial.

If $\operatorname{rk}(G)=1$, that infinite definable centralizer has
rank one. Equation (382.1) would make it finite index, contradicting
the preceding no-finite-index-subgroup argument.\hfill$\square$

This proof excludes rank one for the selected simple groups without
claiming a classification of all rank-one groups. Stronger established
low-rank theorems exist, but are not needed to justify this result.
The obstruction also shows why attempting to construct a finite-rank
counterexample with every proper centralizer finite cannot work.

\paragraph{Definable hulls of abelian subgroups remain abelian.}
For a subset $A\subseteq G$, let $d(A)$ be the intersection of all
definable subgroups containing it. The descending chain argument makes
$d(A)$ definable. If $A$ is an abelian subgroup, then $d(A)$ is abelian.
Indeed, for each $a\in A$, its centralizer is a definable subgroup
containing $A$, so $d(A)\subseteq C_G(a)$. Thus every member of $d(A)$
commutes with every member of $A$. Now fix $x\in d(A)$. The definable
subgroup $C_G(x)$ contains $A$, hence contains $d(A)$. Therefore all
pairs of elements of $d(A)$ commute.

In particular an element of infinite order has an infinite definable
abelian hull of its cyclic subgroup. In a simple nonabelian $G$ this
hull is proper, so has rank strictly between zero and
$\operatorname{rk}(G)$. This yields a controlled lower-rank subgroup
whenever such an element exists. It does not prove that the group has
an infinite-order element, nor that a maximal abelian definable subgroup
is a torus over a field. Those would be additional structure statements.

A related finite-intersection fact says that $C_G(A)$ is the centralizer
of finitely many elements of $A$: it is an intersection of definable
centralizers, to which the same descending chain argument applies.
This controls the syntax of centralizers but does not make their
normalizers or conjugacy geometry identical to those of algebraic tori.

\paragraph{The intended algebraic examples fit the rank bookkeeping.}
Let $K$ be algebraically closed and consider $\operatorname{PGL}_2(K)$.
The basic algebraically-closed-field facts used here are quantifier
elimination and equality of Morley rank with algebraic dimension on
constructible sets. Thus $\operatorname{GL}_2(K)$ has rank four and
its scalar subgroup rank one; the quotient has rank three.
Every determinant has a square root in $K$, so every projective class
has a determinant-one representative. Hence
$\operatorname{PGL}_2(K)\cong\operatorname{PSL}_2(K)$.

One can verify simplicity without treating it as a consequence of the
conjecture. On the projective line, the translations
$u(t)=\left(\begin{smallmatrix}1&t\\0&1\end{smallmatrix}\right)$
form an abelian normal subgroup of the stabilizer of infinity. That
stabilizer is transitive on the remaining affine line, and the whole
group is doubly transitive, hence primitive and faithful.
Translations and their lower-triangular conjugates generate the group
by Gaussian elimination. They are commutators, since
\[
 [\operatorname{diag}(a,a^{-1}),u(t)]=u((a^2-1)t)
\]
and the infinite field contains an $a$ with $a^2\ne1$.
So the projective group is perfect.

For any nontrivial normal subgroup $N$, its orbits are an invariant
partition, hence it is transitive by primitivity and faithfulness.
For each $g$ choose $n\in N$ taking infinity to $g\infty$; then
$g=nh$ with $h$ in the stabilizer. In the quotient by $N$, every
conjugate of the translation subgroup has the same abelian image,
because $h$ normalizes it. Those conjugates generate the quotient,
which is consequently abelian. Perfectness makes it trivial. This
proves abstract simplicity of the infinite algebraic example.

The example illustrates the relevance of rank three, centralizers,
and permutation actions. It is not an argument that an arbitrary
finite-rank simple group has a projective-line action or comes equipped
with a field. Those are precisely the kinds of missing structures a
general identification theorem would need to obtain.

\paragraph{A concrete field reconstruction under an extra splitting hypothesis.}
The sharply two-transitive route in [S1] can be tested in its elementary
split case. Suppose a group acts sharply two-transitively on a set
$\Omega$ and has an abelian normal subgroup $N$ acting regularly.
Regularly means that for each ordered pair of points exactly one
member of $N$ sends the first to the second. Fix $0\in\Omega$ and
identify $\Omega$ with the additive group $N$ through its action on
$0$. Let $H=G_0$.

Each $h\in H$ acts on $N$ by conjugation and hence as an automorphism
of this abelian additive group. Sharp two-transitivity says that $H$
acts regularly on $N\setminus\{0\}$. Choose $1\ne0$. For each
$a\ne0$ let $h_a$ be the unique element with $h_a(1)=a$, and define
\[
 a\cdot b=h_a(b)\quad(a\ne0),\qquad0\cdot b=0.
\]
The additive operation is abelian by hypothesis. Each $h_a$ fixes zero,
so multiplication by a nonzero element sends nonzero elements to
nonzero elements. Since $h_ah_b$ sends $1$ to $a\cdot b$, uniqueness
gives $h_ah_b=h_{a\cdot b}$. This proves associativity on nonzero
factors, and zero factors are immediate. The element $h_1$ is the
identity, and inverses in $H$ give multiplicative inverses.
Finally each $h_a$ is additive, so
\[
 a\cdot(b+c)=a\cdot b+a\cdot c.
 \tag{382.4}
\]
This constructs a one-sided distributive division structure, often
called a near-field; it is not yet automatically a field.

If the stabilizer $H$ is abelian, multiplication is commutative because
$h_ah_b=h_bh_a$. Commutativity and (382.4) then give the other
distributive law too. We have obtained an actual field. The original
group is its affine group: each element uniquely factors as a
translation in $N$ followed by an element of $H$, since its image of
zero first determines the translation. Thus it acts as $x\mapsto ax+b$.
Every step in this reconstruction has its extra hypothesis visible.

\paragraph{Why that reconstruction is a diagnostic, not the classification.}
The affine group just constructed has the proper abelian normal subgroup
of translations and is not simple when $|\Omega|>2$. Its purpose is
to show concretely how a field can emerge from a controlled action and
where extra assumptions enter. It is not the claimed algebraic
realization of every simple group in the question.

First, a sharply two-transitive action need not have a regular abelian
normal subgroup in the general infinite setting. The 2026 source's
nonsplit constructions specifically concern this issue. Second, even
with such a subgroup, regularity of the stabilizer action gives only
the one-sided structure above; commutativity of $H$ supplied the extra
field law. Third, after reconstructing a field, proving that it is
algebraically closed is another substantive model-theoretic theorem.
The calculation does not establish that property by observing that
some associated sets have finite rank.

A proposed counterexample must also put a definable structure of finite
Morley rank on the constructed group or interpret it in such a structure.
Being countable, finitely generated, simple, or sharply transitive does
not establish finite Morley rank. Nor does assigning an integer to the
whole group verify the rank axioms on all parameter-definable subsets
and in saturated elementary extensions.

\paragraph{What a rank induction would still have to prove.}
Proposition 382.1 supplies a proper infinite centralizer, and its
connected component has smaller rank. Even if every simple section
of such lower-rank subgroups had already been identified, it would
remain necessary to understand how these subgroups intersect, normalize
one another, and generate the ambient group. A list of algebraic-looking
proper subgroups is not an isomorphism of the whole group with an
algebraic group.

The connected-component and rank tools also do not imply that conjugation
on a finite collection of tangent spaces gives a faithful finite-dimensional
linear representation: no such tangent spaces have been constructed.
Importing the adjoint representation from algebraic geometry at that
point would assume a principal part of the desired conclusion.

\paragraph{Diagnostics and outcome.}
Finite exact calculations check the displayed commutator and the
near-field reconstruction for affine groups over small prime fields.
They also verify that the reconstructed translation subgroup is normal
and proper, so those examples are not mislabeled simple. Such finite
checks audit identities only; finite groups have rank zero and cannot
test the infinite-rank hypotheses of the conjecture.

\textbf{Outcome: UNRESOLVED.} The attempt proves connectedness and
perfectness restrictions, a complete finite-centralizer obstruction,
the exclusion of rank one in the selected simple case, and an explicit
conditional field-reconstruction mechanism. The first missing step is
to derive enough global geometry and a compatible field or linear
action from the finite-rank axioms for an arbitrary simple group.
Neither the verified rank inequalities nor the source's exotic groups
provide that step.

\subsection{Sources}
\begin{itemize}
\item[S1]K. Tent, From the Cherlin-Zilber Conjecture via sharply 2-transitive groups to the Burnside problem, Abstract (2026).
\url{https://arxiv.org/abs/2606.18207}.
\textit{Formulation source.}
\item[E1]Adrien Deloro, \emph{Tutorial on Groups of finite Morley rank}.
\url{https://personalpages.manchester.ac.uk/staff/mike.prest/DeloroTutorial.pdf}.
\textit{Primary lecture notes inspected 27 September 2026; connectedness,
Morley degree and rank tools are treated as established inputs.}
\item[E2]Joshua Wiscons, \emph{A Short Course on Groups of Finite Morley Rank}.
\url{https://webpages.csus.edu/wiscons/research/GFMR-Minicourse-Notes.pdf}.
\textit{Author-hosted notes inspected 27 September 2026; basic rank and
definability facts are distinguished from an algebraicity theorem.}
\end{itemize}
\end{document}
